\section{Implementation and Reproducibility Details}
\label{app:implementation}

\providecommand{\artifacthash}[4]{\texttt{#1}\allowbreak\texttt{#2}\allowbreak\texttt{#3}\allowbreak\texttt{#4}}
\providecommand{\artifactrevision}[4]{\texttt{#1}\allowbreak\texttt{#2}\allowbreak\texttt{#3}\allowbreak\texttt{#4}}

This appendix records the implementation that produced the completed benchmark
artefacts, rather than reconstructing settings from model names or earlier
configuration drafts.  Encoder settings were read from the provenance sidecars
written during extraction, and estimator settings were read from the completed
run contract and executable entry point.  Paths are reported relative to the
benchmark repository where possible.  Yield values were used directly in their
reported USDA--NASS unit of bushels per acre; no target-unit conversion was
applied.

\subsection{Evaluation cohort and Sentinel-2 composites}

The matched evaluation cohort contains 2,076 county--year observations from 890
counties and covers 2019--2022.  The image corpus used for the completed
extractions contains 427,049 Sentinel-2 patch--timestep files.  These files form
61,007 complete spatial-patch sequences, each with seven timesteps.  The number
of retained patches varies with county area: the completed provenance records a
minimum of 1, a median of 32, and a maximum of 64 patches per county--year.

Each source patch has a spatial size of $256\times256$ pixels at 10-m ground
sampling distance.  The common ten-band order is B02, B03, B04, B05, B06, B07,
B08, B8A, B11, and B12.  The seven composite timestamps are separated by 28
days: 15 April, 13 May, 10 June, 8 July, 5 August, 2 September, and 30
September.  Accordingly, the zero-based month vector supplied to Presto is
$[3,4,5,6,7,8,8]$.  We refer to these inputs as seven 28-day composites
throughout.

Sentinel-2 Level-2A observations were filtered at collection level using
\texttt{CLOUDY\_PIXEL\_PERCENTAGE}$<80\%$.  QA60 and the Scene Classification
Layer were then used to mask cloud, cirrus, cloud-shadow, snow, and ice pixels,
and the remaining observations were combined by a pixel-wise median within
each compositing interval.  A spatial patch was retained only when all seven
timesteps were available.  Oversized arrays were centre-cropped to
$256\times256$; undersized or non-finite arrays caused extraction to terminate.
The completed provenance records zero undersized or incomplete spatial patches.
The saved patch artefacts do not retain the complete list of source Sentinel-2
scene identifiers, so scene-level reconstruction is not possible from the
benchmark directory alone.

\subsection{Exact encoder variants and model artefacts}

Table~\ref{tab:encoder_checkpoints_exact} identifies the implementation and
checkpoint recorded by each completed embedding run.  The SHA-256 values below
the table were computed from the model artefacts copied into the benchmark
directory.  They make the local copies verifiable, although the extraction
sidecars themselves recorded paths and model identifiers rather than hashes.

\begin{table*}[t]
\centering
\caption{Exact encoder implementation and model artefact associated with the
completed embedding files.  ``Dim.'' is the per-row embedding dimension before
county aggregation.}
\label{tab:encoder_checkpoints_exact}
\scriptsize
\setlength{\tabcolsep}{4pt}
\renewcommand{\arraystretch}{1.12}
\begin{tabular}{p{2.2cm}p{3.2cm}p{8.4cm}c}
\toprule
\textbf{Representation} & \textbf{Implementation identifier} &
\textbf{Checkpoint or precomputed product used by the completed run} &
\textbf{Dim.} \\
\midrule
Clay & \path{clay_v1_5} & Official \path{clay-v1.5.ckpt}, loaded through the
final encoder in the Clay model repository; the accompanying official
\path{configs/metadata.yaml} supplies the Sentinel-2 statistics, wavelengths,
and 10-m ground sampling distance. & 1024 \\

Prithvi-EO-2.0-300M-TL & \path{prithvi_eo_v2_300_tl} &
\path{ibm-nasa-geospatial/Prithvi-EO-2.0-300M-TL}, loaded through TerraTorch
1.2.11.  The cached Hugging Face snapshot revision is
\path{63adbd39c271da4c42f447e69b1a7c91a338cdc9}. & 1024 \\

TerraMind, 6 bands & \path{terramind_v1_base_s2_6_prithvi} &
\path{ibm-esa-geospatial/TerraMind-1.0-base}, file
\path{TerraMind_v1_base.pt}, loaded through TerraTorch 1.2.11.  The cached
snapshot revision is \path{fb96c70d0a5f68dcc44030b89cbfd8ec3fb0c67a}. & 768 \\

TerraMind, 10 bands & \path{terramind_v1_base_s2_10_zero_pad} & Same TerraMind
v1 base checkpoint and cached revision as the six-band experiment; only the
spectral adapter differs. & 768 \\

Presto & \path{presto_s2} & \path{nasaharvest/presto} revision
\artifactrevision{ba88a3f9d3}{06ea30b90f}{e2f1284d98}{f423487cfa}, using the bundled
\path{data/default_model.pt}. & 128 \\

AlphaEarth V1 & \path{GOOGLE/SATELLITE_EMBEDDING/V1/ANNUAL} & Google Satellite
Embedding V1 annual product, retrieved through Earth Engine and reduced there
to county-level means.  No local encoder checkpoint was used.  The exported
vectors contain the 64 axes A00--A63. & 64 \\
\bottomrule
\end{tabular}
\end{table*}

The copied Clay checkpoint has SHA-256
\artifacthash{21432069250b9b3f}{9a65ffd0071c5ad5}{6b793247285ab060}{4edf7f531d4798d0}.
The copied Presto checkpoint has SHA-256
\artifacthash{559c9d76f525bdf2}{dd3774da8b1e7f15}{10769b73baacac12}{e9e4d8d5c67fbe65}.
The cached TerraMind v1 base file has SHA-256
\artifacthash{83c3a0938067c838}{67a46e564443c2fa}{38383bf4f966d931}{b11cb025b847d7ec},
and the cached Prithvi-EO-2.0-300M-TL file has SHA-256
\artifacthash{3629cedfbb350faaf}{cb0dac902ae0d3c9}{27e25ce8d9e0024a}{a1276ec66956ddb}.

The repository also contains
\path{models/Prithvi-EO-2.0-300M/Prithvi_EO_V2_300M.pt} (SHA-256
\artifacthash{faab2c8b776f0723}{f93431747603a7ec}{586a54fc45473613}{b4322c29752486e7}).
This is the non-TL 300M checkpoint, but it is \emph{not} the checkpoint named
by the saved Prithvi provenance.  We therefore report the completed experiment
as Prithvi-EO-2.0-300M-TL.  The saved file contains 427,049 independently
encoded, 1024-dimensional patch--timestep vectors; using the non-TL checkpoint
would require regenerating the embeddings and all downstream results.

For this study, the AlphaEarth vectors were retrieved from the Earth Engine
collection
\url{https://developers.google.com/earth-engine/datasets/catalog/GOOGLE_SATELLITE_EMBEDDING_V1_ANNUAL}.
The V1 annual product was used at the time of the paper.  The 64 embedding axes
A00--A63 were masked to corn pixels with the USDA Cropland Data Layer and
averaged over each county in Earth Engine, producing one vector per
county--year.  The resulting county table, rather than a local AlphaEarth model
checkpoint or patch-level embedding file, is the input to the benchmark.

\subsection{Embedding extraction}

\paragraph{Clay.}
All ten Sentinel-2 bands were used at $256\times256$ pixels.  Source digital
numbers were converted to reflectance when required, and the checkpoint's
Sentinel-2 means and standard deviations were converted to the same units
before standardisation.  The extractor used the official cyclic week/hour and
latitude/longitude encodings and wavelengths in micrometres.  The output is the
1024-dimensional CLS token from the final encoder layer, one vector per patch
and timestep.  The completed file contains 427,049 vectors.

\paragraph{Prithvi-EO-2.0-300M-TL.}
The source image was first checked against the $256\times256$ benchmark
footprint and then centre-cropped to the native $224\times224$ model input,
without interpolation or padding.  Six bands were selected in the order B02,
B03, B04, B8A, B11, and B12 and supplied in digital-number units.  The official
Prithvi V2 means $(1087,1342,1433,2734,1958,1363)$ and standard deviations
$(2248,2179,2178,1850,1242,1049)$ were applied once.  Each composite was
encoded independently as a single frame, with its year, one-based day of year,
latitude, and longitude.  The 196 non-CLS tokens from the final layer were
spatially averaged to obtain one 1024-dimensional vector per patch and
timestep, giving 427,049 output rows.

\paragraph{TerraMind.}
Both experiments used TerraMind v1 base and a $224\times224$ centre crop.  The
input-unit detector converted source values to reflectance and clipped them to
$[0,1]$.  The six-band experiment passed B02, B03, B04, B8A, B11, and B12
directly through the model's band-selection interface, with no zero padding.
The ten-band experiment mapped B02, B03, B04, B05, B06, B07, B08, B8A, B11,
and B12 into TerraMind's 12-channel Sentinel-2 adapter and inserted zeros only
for unavailable B01 and B09.  All 196 final-layer spatial tokens were averaged,
producing one 768-dimensional vector per patch and timestep.  Each experiment
contains 427,049 rows.

\paragraph{Presto.}
Presto was the only frozen encoder that received the complete seven-step
sequence in one forward pass.  After the $256\times256$ crop, each band was
spatially averaged over the patch.  Source units were detected per file and
converted to raw digital numbers before calling
\path{presto.construct_single_presto_input}; the official helper was responsible
for normalisation, masking, and NDVI construction.  The completed experiment
used the spatial-mean mode, not pixel sampling, and did not provide ERA5-Land to
the encoder.  The official evaluation output was globally pooled to one
128-dimensional vector per complete spatial-patch sequence, yielding 61,007
rows.

\paragraph{AlphaEarth and handcrafted features.}
Earth Engine performed the spatial reduction for AlphaEarth.  The imported
table therefore already contained one county--year row with the 64 columns
\path{mean_A00}--\path{mean_A63}; these values were used without further
pooling or normalisation.  The Sentinel-2 baseline similarly used 21
precomputed county--year values: EVI, LAI, and fPAR for each of the seven
28-day composites.

\subsection{County aggregation used by the completed main table}

For a county--year with embedding rows $\{\mathbf{z}_i\}_{i=1}^{n}$, the
completed main-table entry point constructed
\begin{equation}
  \mathbf{x}=
  \left[
    \operatorname{mean}_{i}(\mathbf{z}_i),\;
    \operatorname{std}_{i}(\mathbf{z}_i;\,\mathrm{ddof}=0)
  \right].
\end{equation}
For Clay, Prithvi, and TerraMind, $i$ indexes patch--timestep rows; the completed
main table therefore pooled space and time jointly and did not retain temporal
order.  Their final widths were 2048, 2048, and 1536, respectively.  For
Presto, each row already represents a complete temporal sequence, so the same
operation pools only across spatial patches and produces 256 features.
AlphaEarth and the 21 Sentinel-2 indices were already county--year vectors and
were used as supplied.

This operation is specific to the completed main-table results.  The temporal
ablation and the climate-free LOYO implementation use the intended two-stage
operation: population mean and standard deviation across the same complete
patches at each timestep, followed by the experiment-specific temporal
operation.  A rerun that replaces the completed main table with this two-stage
operation is a new experiment and should not be presented as the source of the
saved scores below.

More precisely, let $\mathbf{z}_{p,t}\in\mathbb{R}^{D}$ denote the embedding
for patch $p$ at timestep $t$.  Two-stage pooling first constructs
\begin{equation}
  \mathbf{s}_{t}=
  \left[
    \operatorname{mean}_{p}(\mathbf{z}_{p,t}),\;
    \operatorname{std}_{p}(\mathbf{z}_{p,t};\,\mathrm{ddof}=0)
  \right]\in\mathbb{R}^{2D},
\end{equation}
and only then applies temporal mean pooling, concatenation, or the learned
one-dimensional convolution to the sequence
$\{\mathbf{s}_{t}\}_{t=1}^{7}$.  Joint pooling instead treats every
patch--timestep row as a member of one unordered set and computes a single
mean and standard deviation across both axes.  It therefore removes temporal
order immediately and combines within-timestep spatial variation with changes
in the timestep means.  The mean components agree when every timestep has the
same number of patches, but the standard-deviation components generally do
not.

\subsection{Splits, fitting, and regression heads}

The primary comparison uses five county-grouped folds generated with
\texttt{GroupKFold} on county FIPS.  All years from a county remain in the same
outer fold.  For outer fold $f$, fold $f$ is the test partition, fold
$(f+1)\bmod 5$ is the validation partition, and the remaining three folds are
used for training.  Every year from 2019 to 2022 occurs in the training,
validation, and test partitions of every outer fold.  Fixed nonlinear models
were fitted on training plus validation data and evaluated once on the outer
test fold.  Ridge used the validation partition to select its penalty before
being refitted on training plus validation data.

\begin{table*}[t]
\centering
\caption{Regression settings used by the completed five-fold main-table run.}
\label{tab:regressor_settings_exact}
\scriptsize
\setlength{\tabcolsep}{4pt}
\renewcommand{\arraystretch}{1.12}
\begin{tabular}{p{2.1cm}p{13.6cm}}
\toprule
\textbf{Head} & \textbf{Exact fitting settings} \\
\midrule
Ridge & Training-only \texttt{StandardScaler}; $\alpha\in
\{0.01,0.1,1,10,100\}$ selected by validation RMSE; scaler and Ridge refitted
on training plus validation data.  Ridge was run once because it is
deterministic. \\

Random Forest & 600 trees; squared-error criterion; unlimited depth;
\texttt{min\_samples\_leaf=2}; \texttt{max\_features=1.0}; bootstrap enabled;
\texttt{random\_state}$\in\{0,1,2\}$. \\

XGBoost & 600 estimators; learning rate 0.03; depth 6;
\texttt{min\_child\_weight=2}; row and column subsampling 0.8; $L_1=0$;
$L_2=1$; squared-error objective; histogram tree method;
\texttt{random\_state}$\in\{0,1,2\}$. \\

EBM & \texttt{interpret} 0.7.8
\texttt{ExplainableBoostingRegressor}; \texttt{interactions=0};
\texttt{max\_bins=1024}; \texttt{max\_interaction\_bins=64};
\texttt{validation\_size=0.15}; 14 outer bags; no inner bags; learning rate
0.04; greedy ratio 10; non-cyclic boosting; 50,000 maximum rounds; 100-round
early stopping with tolerance $10^{-5}$; 500 main-effect and 100 interaction
smoothing rounds; \texttt{min\_samples\_leaf=4}; minimum Hessian 0; $L_1=0$;
$L_2=0$; maximum delta step 0; gain scale 5; two leaves; missing values treated
as a separate category; RMSE objective;
\texttt{random\_state}$\in\{0,1,2\}$; four CPU workers. \\
\bottomrule
\end{tabular}
\end{table*}

For the stochastic heads, predictions and metrics were averaged across seeds
0, 1, and 2 within each outer fold.  Reported central values and standard
deviations are then the mean and population standard deviation across the five
outer folds.  Metrics are county-level $R^2$, RMSE, and MAE, with RMSE and MAE
reported in bushels per acre.

The saved run contract records the cohort, representations, heads, seeds, and
split path, but not the package versions or worker count.  The versions above
come from the repository's recorded regression environment, and four workers
is the documented entry-point default.  Retaining an environment lock file and
the full command line with future runs would remove this remaining ambiguity.

\subsection{Completed EBM run}

The EBM run is complete for all representations in the saved main-table
artefact.  The implementation used the \texttt{interpret} 0.7.8 defaults listed
in Table~\ref{tab:regressor_settings_exact}, except that pairwise interactions
were disabled.  This produced an additive EBM and avoided the library default
\texttt{interactions='5x'}, which would request 10,240 pairwise terms for a
2048-dimensional representation.  Table~\ref{tab:ebm_completed} reports the
saved five-fold results.  The Sentinel-2+Daymet row is an auxiliary climate
late-fusion result and is not part of the climate-free encoder comparison.

\begin{table}[t]
\centering
\caption{Completed EBM results.  $R^2$ and RMSE are mean $\pm$ population
standard deviation across five county-grouped outer folds after seed averaging;
MAE is the corresponding fold mean.}
\label{tab:ebm_completed}
\scriptsize
\setlength{\tabcolsep}{3.2pt}
\renewcommand{\arraystretch}{1.08}
\begin{tabular}{lccc}
\toprule
\textbf{Representation} & $\mathbf{R^2}$ & \textbf{RMSE} & \textbf{MAE} \\
\midrule
AlphaEarth & $0.8260\pm0.0312$ & $14.3080\pm1.1126$ & 10.2626 \\
S2 indices & $0.7589\pm0.0484$ & $16.8122\pm1.0604$ & 12.7402 \\
S2 + Daymet$^{\dagger}$ & $0.8595\pm0.0172$ & $12.8879\pm0.4621$ & 9.8236 \\
Clay & $0.8448\pm0.0382$ & $13.4540\pm1.1379$ & 10.0715 \\
Prithvi-EO-2.0-300M-TL & $0.8177\pm0.0402$ & $14.6060\pm1.0789$ & 10.8846 \\
TerraMind, 6 bands & $0.8067\pm0.0449$ & $15.0317\pm1.2453$ & 11.3173 \\
TerraMind, 10 bands & $0.8169\pm0.0483$ & $14.5964\pm1.3135$ & 10.8417 \\
Presto & $0.8000\pm0.0493$ & $15.2783\pm1.3591$ & 11.6661 \\
\bottomrule
\multicolumn{4}{l}{\scriptsize $^{\dagger}$Auxiliary county-level climate late fusion.}
\end{tabular}
\end{table}

\subsection{Climate late fusion and temporal readout ablation}

Daymet was used only after county aggregation.  The feature order is day
length, precipitation, shortwave radiation, maximum temperature, and minimum
temperature, with seven 28-day values for each variable, giving 35 climate
features.  Concatenating these with the 21 Sentinel-2 EVI/LAI/fPAR values gives
the 56-dimensional Sentinel-2+Daymet vector.  Daymet was not passed through
Clay, Prithvi, TerraMind, or Presto.  The completed main-table directory does
not contain a native Presto+ERA5-Land result.

The completed temporal-ablation artefacts cover all five outer folds and seeds
0, 1, and 2 for Clay, Prithvi, and six-band TerraMind.  Complete patch identity
is matched across encoders.  At each timestep, patch embeddings are represented
by their population mean and population standard deviation, producing a
seven-step county sequence.  The \texttt{mean} and \texttt{concat} strategies
use a learned 256-unit projection; \texttt{conv1d} uses two temporal
convolutions with kernel size 3 and 128 then 256 channels, followed by temporal
global mean pooling.  All strategies share a LayerNorm--128-unit MLP head with
GELU and dropout 0.2.  Inputs are standardised per timestep and feature, and
targets are standardised, using the training partition only.  Training uses
AdamW, learning rate $10^{-3}$, weight decay $10^{-4}$, batch size 32, MSE on
the standardised target, at most 300 epochs, and early stopping after 30 epochs
without improvement in validation county RMSE.  Mixed precision was disabled
in the completed run.

\subsection{Supervised 3D-ConvLSTM configuration}

The supervised Sentinel-2 reference uses the same seven-composite,
ten-channel, $256\times256$ input contract and the same county-grouped outer
folds.  Its 3-D convolutional stem has channels $[24,48,64]$, kernel
$3\times3\times3$, spatial stride 2 at each stage, group normalisation with
eight groups, GELU, and dropout 0.15.  A single ConvLSTM layer with 64 hidden
channels and a $3\times3$ kernel carries spatial hidden and cell states through
the sequence.  The final hidden state is spatially pooled and passed through a
64-unit regression head.

Training uses a learning rate of $3\times10^{-4}$, weight decay $10^{-5}$,
batch size 1, eight sampled patches per county--year, patch chunk size 2,
four-step gradient accumulation, gradient clipping at 1.0, mixed precision, a
maximum of 50 epochs, and early stopping patience 8.  Band and target
normalisation statistics are fitted on training county--years only.  The
validation fold selects the best epoch by RMSE and the outer test fold is
evaluated once.  This appendix records the executable configuration; no
completed supervised result is claimed here.

\subsection{Software, hardware, and saved audit artefacts}

Frozen embedding extraction ran on the JUPITER booster partition using one
NVIDIA GH200 GPU under Python 3.13.5 and PyTorch 2.13.0+cu130.  Prithvi and
TerraMind used TerraTorch 1.2.11.  Presto used its separate environment and the
repository revision shown in Table~\ref{tab:encoder_checkpoints_exact}.  The
recorded local regression environment was Python 3.13.13, NumPy 2.5.2, pandas
3.0.5, scikit-learn 1.9.0, XGBoost 3.4.1, and \texttt{interpret} 0.7.8.  The
completed temporal ablation records PyTorch 2.12.0, NumPy 2.4.6, and pandas
3.0.3.

The checkpoint cache is retained under \path{hf_cache/}, including the exact
TerraMind and Prithvi snapshot revisions listed above.  The copied model
artefacts are under \path{clay/} and \path{models/}.  Each completed embedding
has a sidecar named \path{outputs/embeddings/*.parquet.provenance.json}, which
records the model identifier, checkpoint, source and output counts, band and
shape contracts, pooling operation, and exclusion counts.  The completed
regression results and their cohort/head contract are stored in
\path{outputs/main_table_covered/main_table.csv} and
\path{outputs/main_table_covered/run_contract.json}.  The temporal-ablation
folds additionally save data contracts, checkpoints, predictions, per-seed
metrics, and fold summaries under \path{outputs/temporal_ablation_covered/}.
